\setcounter{section}{4}

  \zwischenueberschrift{Die Weingartenabbildung}

Es sei

\mavergleichskette
{\vergleichskette
{ Y
}
{ \subseteq }{ \R^n
}
{  }{
}
{    }{
}
{   }{
}
}
{}{}{}
eine
\definitionsverweis {differenzierbare Hyperfläche}{}{.}
Zu einem Punkt

\mavergleichskette
{\vergleichskette
{ P
}
{ \in }{ Y
}
{  }{
}
{    }{
}
{   }{
}
}
{}{}{}
definiert man eine lineare Abbildung
\maabbdisp {L_P} { T_PY } { T_PY
} {}
auf dem Tangentialraum zu $Y$. Es sei $N$ ein differenzierbares Einheitsnormalenfeld auf $Y$, das auf einer offenen Umgebung von $Y$ definiert sei. Die wesentliche Idee ist, einen Tangentialvektor

\mavergleichskette
{\vergleichskette
{ v
}
{ \in }{ T_PY
}
{  }{
}
{    }{
}
{   }{
}
}
{}{}{}
durch eine differenzierbare Kurve
\maabbdisp {\gamma} { I } { Y
} {}
zu parametrisieren, dabei ist also

\mavergleichskettedisp
{\vergleichskette
{ \gamma'(0)
}
{ =} { v
}
{ } {
}
{ } {
}
{ } {
}
}
{}{}{.}
Das Einheitsnormalenfeld definiert dann die Abbildung
\maabbdisp {N \circ \gamma} { I } { \R^n
} {}
längs $I$. Die infinitesimale Änderung des Einheitsnormalenfeldes längs $\gamma$ wird durch den Limes

\mathdisp {\operatorname{lim}_{   t  \rightarrow  0   } \,  { \frac{   N(\gamma(t))-N(P)  }{  t  } }} {  }

gemessen, falls dieser existiert. Nach
[[Richtungsableitung/K/Lineare Realisierung/Differenzierbare Kurve/Fakt|Lemma 43.4 (Analysis (Osnabrück 2021-2023))]]
ist dieser Limes gleich der Richtungsableitung $\left(  D_{ v }  N   \right) \left(   P   \right)$ und wiederum gleich
\mathl{\left(  D N   \right)_{ P } \left(   v   \right)}{,} dem
\definitionsverweis {totalen Differential}{}{}
ausgewertet am Vektor $v$. Es wird sich herausstellen, dass diese Zuordnung
\maabbeledisp {} { T_PY } { \R^n
} {v } { \left(  D_{v}  N   \right) \left(   P   \right)
} {,}
in $T_PY$ landet.

\inputdefinition
{}
{

Es sei

\mavergleichskette
{\vergleichskette
{ W
}
{ \subseteq }{ \R^n
}
{  }{
}
{    }{
}
{   }{
}
}
{}{}{}
\definitionsverweis {offen}{}{,}
\maabb {h} { W } { \R
} {}
eine zweifach
\definitionsverweis {stetig differenzierbare Funktion}{}{}
und

\mavergleichskette
{\vergleichskette
{ Y
}
{ = }{ h^{-1}(c)
}
{  }{
}
{    }{
}
{   }{
}
}
{}{}{}
die
\definitionsverweis {Faser}{}{}
zu

\mavergleichskette
{\vergleichskette
{ c
}
{ \in }{ \R
}
{  }{
}
{    }{
}
{   }{
}
}
{}{}{,}
wobei $h$ in jedem Punkt von $Y$
\definitionsverweis {regulär}{}{}
sei. Es sei $N$ ein
\definitionsverweis {Einheitsnormalenfeld}{}{}
und sei

\mavergleichskette
{\vergleichskette
{ P
}
{ \in }{ Y
}
{  }{
}
{    }{
}
{   }{
}
}
{}{}{.}
Dann nennt man
\maabbeledisp {L_P} {T_P Y} { T_PY
} { v } { - \left(  D_{v}  N   \right) \left(   P  \right)
} {,}
die
\definitionswort {Weingartenabbildung}{}
in $T_PY$.

}

Man beachte, dass die Weingartenabbildung vom Einheitsnormalenfeld abhängt, auch wenn dies nicht immer explizit gesagt wird. Wenn $Y$ als Faser zu $h$ gegeben ist, so nimmt man in der Regel das zugehörige normierte Gradientenfeld zu $h$.

__NOEDITSECTION__

\inputfaktbeweis
{Hyperfläche/Weingartenabbildung/Tangentialraum/Fakt}
{Lemma}
{}
{

\faktsituation {Es sei

\mavergleichskette
{\vergleichskette
{ W
}
{ \subseteq }{ \R^n
}
{  }{
}
{    }{
}
{   }{
}
}
{}{}{}
\definitionsverweis {offen}{}{,}
\maabb {h} { W } { \R
} {}
eine zweifach
\definitionsverweis {stetig differenzierbare Funktion}{}{}
und

\mavergleichskette
{\vergleichskette
{ Y
}
{ = }{ h^{-1}(c)
}
{  }{
}
{    }{
}
{   }{
}
}
{}{}{}
die
\definitionsverweis {Faser}{}{}
zu

\mavergleichskette
{\vergleichskette
{ c
}
{ \in }{ \R
}
{  }{
}
{    }{
}
{   }{
}
}
{}{}{,}
wobei $h$ in jedem Punkt von $Y$
\definitionsverweis {regulär}{}{}
sei. Es sei

\mavergleichskette
{\vergleichskette
{ P
}
{ \in }{ Y
}
{  }{
}
{    }{
}
{   }{
}
}
{}{}{.}}
\faktfolgerung {Dann ist die
\definitionsverweis {Weingartenabbildung}{}{}
$L_P$ ein linearer Endomorphismus des
\definitionsverweis {Tangentialraumes}{}{}
$T_pY$.}
\faktzusatz {}
\faktzusatz {}

}
{

Es ist

\mavergleichskette
{\vergleichskette
{ N(P)
}
{ = }{ { \frac{
\operatorname{Grad} \,  h  (  P )  }{  \Vert {
\operatorname{Grad} \,  h  (  P ) } \Vert  } }
}
{  }{
}
{    }{
}
{   }{
}
}
{}{}{}
ein stetig differenzierbares Vektorfeld, das auf einer offenen Umgebung

\mavergleichskette
{\vergleichskette
{ Y
}
{ \subseteq }{ U
}
{  }{
}
{    }{
}
{   }{
}
}
{}{}{}
definiert ist. Daher ist gemäß
[[Differenzierbarkeit/R/Totale Differenzierbarkeit impliziert richtungsweise Differenzierbarkeit/Fakt|Proposition 46.1 (Analysis (Osnabrück 2021-2023))]]

\mavergleichskettedisp
{\vergleichskette
{ \left(  D_{v}  N   \right) \left(   P  \right)
}
{ =} { \left(  D N   \right)_{ P } \left(  v  \right)
}
{ } {
}
{ } {
}
{ } {
}
}
{}{}{}
linear in der Richtung $v$. Wegen der Einheitsnormalenbedingung ist

\mavergleichskette
{\vergleichskette
{ \left\langle  N(P)  ,  N(P)   \right\rangle
}
{ = }{ 1
}
{  }{
}
{    }{
}
{   }{
}
}
{}{}{}
für alle

\mavergleichskette
{\vergleichskette
{ P
}
{ \in }{ U
}
{  }{
}
{    }{
}
{   }{
}
}
{}{}{}
und daher ist unter Verwendung von
[[Skalarprodukt von Abbildungen/Richtungsableitung/Aufgabe|Aufgabe 43.14 (Analysis (Osnabrück 2021-2023))]]

\mavergleichskettedisp
{\vergleichskette
{ 0
}
{ =} { \left(  D_{v}  \left\langle  N  ,  N  \right\rangle   \right) \left(   P  \right)
}
{ =} { 2 \left\langle  N(P)  ,  \left(  D_{v}  N   \right) \left(   P  \right)   \right\rangle
}
{ } {
}
{ } {
}
}
{}{}{.}
Daher steht
\mathl{\left(  D_{v}  N   \right) \left(   P  \right)}{} senkrecht auf $N(P)$ und gehört bei

\mavergleichskette
{\vergleichskette
{ P
}
{ \in }{ Y
}
{  }{
}
{    }{
}
{   }{
}
}
{}{}{}
zum Tangentialraum $T_PY$.

}

Die Weingartenabbildung $L_P$ ist also die negierte Einschränkung des totalen Differentials des Einheitsnormalenfeldes auf den Tangentialraum $T_PY$ in den Tangentialraum $T_PY$. Wenn man das totale Differential über die
\definitionsverweis {Jacobi-Matrix}{}{}
von $-N$ ausrechnet, so muss man deren Wirkungsweise auf einer Basis des Tangentialraumes bestimmen, um eine Matrixdarstellung der Weingartenabbildung zu erhalten.

__NOEDITSECTION__

\inputfaktbeweis
{Ebene Kurve/Weingartenabbildung/Krümmung/Fakt}
{Lemma}
{}
{

\faktsituation {Es sei

\mavergleichskette
{\vergleichskette
{ W
}
{ \subseteq }{ \R^2
}
{  }{
}
{    }{
}
{   }{
}
}
{}{}{}
\definitionsverweis {offen}{}{,}
\maabb {h} { W } { \R
} {}
eine zweifach
\definitionsverweis {stetig differenzierbare Funktion}{}{}
und

\mavergleichskette
{\vergleichskette
{ C
}
{ = }{ h^{-1}(c)
}
{  }{
}
{    }{
}
{   }{
}
}
{}{}{}
die Faser zu

\mavergleichskette
{\vergleichskette
{ c
}
{ \in }{ \R
}
{  }{
}
{    }{
}
{   }{
}
}
{}{}{,}
wobei $h$ in jedem Punkt von $C$
\definitionsverweis {regulär}{}{}
sei und es sei

\mavergleichskette
{\vergleichskette
{ P
}
{ \in }{ C
}
{  }{
}
{    }{
}
{   }{
}
}
{}{}{.}}
\faktfolgerung {Dann ist die
\definitionsverweis {Weingartenabbildung}{}{}
in $P$ die Multiplikation mit der
\definitionsverweis {Krümmung}{}{}
$\kappa(P)$ von $C$ in $P$.}
\faktzusatz {}
\faktzusatz {}

}
{

Dies ist eine Umformulierung von
[[Ebene differenzierbare Kurve/Faser/Krümmung/Fakt|Lemma 3.11]].

}

In der vorstehenden Aussage wurde nicht explizit auf die Orientierung Bezug genommen, dies muss man sich dazudenken.

\inputbeispiel{}
{

Es sei

\mavergleichskettedisp
{\vergleichskette
{ h(x,y,z)
}
{ =} { x^2+y^2+z^2
}
{ } {
}
{ } {
}
{ } {
}
}
{}{}{}
und wir betrachten die Faser $Y$ zu $h$ über $r^2$, also die Kugeloberfläche zum Radius

\mavergleichskette
{\vergleichskette
{ r
}
{ > }{ 0
}
{  }{
}
{    }{
}
{   }{
}
}
{}{}{}
mit dem Ursprung als Mittelpunkt. Es sei

\mavergleichskette
{\vergleichskette
{ P
}
{ = }{ \left(  x_0   , \,   y_0  , \,  z_0                 \right)
}
{ \in }{ Y
}
{    }{
}
{   }{
}
}
{}{}{.}
Durch eine Isometrie kann man diesen Punkt nach

\mavergleichskette
{\vergleichskette
{ Q
}
{ = }{ \left(  r   , \,   0  , \,   0                 \right)
}
{  }{
}
{    }{
}
{   }{
}
}
{}{}{}
transformieren, was die
\definitionsverweis {Weingartenabbildung}{}{}
nicht ändert. Eine Basis des Tangentialraumes ist dann
\mathkor {} {\begin{pmatrix}  0  \\ 1 \\  0   \end{pmatrix}} {und} {\begin{pmatrix}  0  \\ 0\\  1   \end{pmatrix}} {.}
Das nach innen gerichtete Einheitsnormalenfeld $N$  ist
\mathl{- { \frac{   1  }{  2r } } \begin{pmatrix}  2x \\ 2y\\  2z   \end{pmatrix}}{} und daher ist zu

\mavergleichskettedisp
{\vergleichskette
{ v
}
{ =} { \begin{pmatrix}  0  \\ b \\ c   \end{pmatrix}
}
{ } {
}
{ } {
}
{ } {
}
}
{}{}{}

\mavergleichskettedisp
{\vergleichskette
{ \left(  D_{v}  N   \right) \left(   Q   \right)
}
{ =} { b { \frac{   \partial   N   }{  \partial   y  } } (Q) +  c { \frac{   \partial   N   }{  \partial  z  } }  (Q)
}
{ =} { -b { \frac{   1  }{  2r } } \begin{pmatrix}  0  \\ 2 \\  0   \end{pmatrix} - c { \frac{   1  }{  2r } } \begin{pmatrix}  0  \\ 0\\  2   \end{pmatrix}
}
{ =} { - { \frac{   1  }{ r } } \begin{pmatrix}  0  \\ b \\ c   \end{pmatrix}
}
{ } {
}
}
{}{}{.}
Daher ist die Weingartenabbildung Multiplikation mit dem Kehrwert ${ \frac{   1  }{ r } }$ des Radius.

}

\inputbeispiel{}
{

Wir knüpfen an
[[Einschaliges Hyperboloid/Tangentenebene/Normalenvektor/Beispiel|Beispiel 2.5]]
an. Die Jacobi-Matrix des Einheitsnormalenfeldes

\mavergleichskettedisphandlinks
{\vergleichskettedisphandlinks
{ N( \begin{pmatrix}  x  \\ y \\ z   \end{pmatrix} )
}
{ =} { \left(  x^2+y^2+z^2  \right)^{ - { \frac{   1  }{  2 } } } \begin{pmatrix}  x  \\ y \\  -z   \end{pmatrix}
}
{ } {
}
{ } {
}
{ } {
}
}
{}{}{}
ist

\mathdisp {\left(  x^2+y^2+z^2  \right)^{- { \frac{   3  }{  2 } } } \begin{pmatrix}  y^2+z^2  &   -xy &  -xz \\  -xy &  x^2+z^2 &  -yz \\ xz &  yz &  -x^2-y^2  \end{pmatrix}} { . }

Die
\definitionsverweis {Weingartenabbildung}{}{}
ist die negierte Einschränkung dieser Abbildung auf den Tangentialraum an die Fläche, wobei
[[Hyperfläche/Weingartenabbildung/Tangentialraum/Fakt|Lemma 4.2]]
sicherstellt, dass wir wieder im Tangentialraum landen. Wir setzen

\mavergleichskette
{\vergleichskette
{ x
}
{ \neq }{ 0
}
{  }{
}
{    }{
}
{   }{
}
}
{}{}{}
voraus und arbeiten mit der Basis
\mathkor {} {\begin{pmatrix}  y  \\ -x\\  0   \end{pmatrix}} {und} {\begin{pmatrix}  z  \\ 0 \\  x   \end{pmatrix}} {}
des Tangentialraumes. Es ist

\mavergleichskettealignhandlinks
{\vergleichskettealignhandlinks
{ \begin{pmatrix}  y^2+z^2  &   -xy &  -xz \\  -xy &  x^2+z^2 &  -yz \\ xz &  yz &  -x^2-y^2  \end{pmatrix} \begin{pmatrix}  y  \\ -x\\  0   \end{pmatrix}
}
{ =} { \begin{pmatrix}  y \left(  y^2+z^2  \right) +x^2y \\ -xy^2 -x \left(  x^2+z^2  \right) \\  0   \end{pmatrix}
}
{ =} { \left(  x^2+y^2+z^2  \right) \begin{pmatrix}  y  \\ -x\\  0   \end{pmatrix}
}
{ } {
}
{ } {
}
}
{}
{}{,}
sodass wir unmittelbar den Eigenvektor $\begin{pmatrix}  y  \\ -x\\  0   \end{pmatrix}$ mit dem Eigenwert
\mathl{- \left(  x^2+y^2+z^2  \right)^{ - { \frac{   1  }{  2 } } }}{} gefunden haben. Ferner ist

\mavergleichskettedisphandlinks
{\vergleichskettedisphandlinks
{ \begin{pmatrix}  y^2+z^2  &   -xy &  -xz \\  -xy &  x^2+z^2 &  -yz \\ xz &  yz &  -x^2-y^2  \end{pmatrix} \begin{pmatrix}  z  \\ 0 \\  x   \end{pmatrix}
}
{ =} { \begin{pmatrix} z \left(  y^2+z^2  \right) -x^2z \\ -2xyz\\  xz^2-x \left(  x^2+y^2  \right)   \end{pmatrix}
}
{ =} { 2yz \begin{pmatrix}  y  \\ -x\\  0   \end{pmatrix} + \left(  z^2-x^2-y^2  \right) \begin{pmatrix}  z  \\ 0 \\  x   \end{pmatrix}
}
{ =} { 2yz \begin{pmatrix}  y  \\ -x\\  0   \end{pmatrix}- \begin{pmatrix}  z  \\ 0 \\  x   \end{pmatrix}
}
{ } {
}
}
{}{}{.}
Somit wird die Weingartenabbildung bezüglich dieser Basis durch die Matrix

\mathdisp {\begin{pmatrix}  - \left(  x^2+y^2+z^2  \right)^{ - { \frac{   1  }{  2 } } }   &   - 2yz \left(  x^2+y^2+z^2  \right)^{ - { \frac{   3  }{  2 } } }    \\  0  &  \left(  x^2+y^2+z^2  \right)^{ - { \frac{   3  }{  2 } } }   \end{pmatrix}} {  }

beschrieben. Einen zweiten Eigenvektor im Tangentialraum erhält man
\zusatzklammer {in Hinblick auf
[[Hyperfläche/Weingartenabbildung/Selbstadjungiert/Eigenwerte/Fakt|Korollar 4.8]]} {} {}
am einfachsten, wenn man zum ersten Eigenvektor und zum Normalenvektor einen senkrechten Vektor bestimmt. Dies ergibt den Vektor
\mathl{\begin{pmatrix}  xz \\ yz\\  x^2+y^2   \end{pmatrix}}{,} und in der Tat ist
\zusatzklammer {ohne den skalaren Vorfaktor} {} {}

\mavergleichskettealignhandlinks
{\vergleichskettealignhandlinks
{ \begin{pmatrix}  y^2+z^2  &   -xy &  -xz \\  -xy &  x^2+z^2 &  -yz \\ xz &  yz &  -x^2-y^2  \end{pmatrix} \begin{pmatrix}  xz \\ yz\\  x^2+y^2   \end{pmatrix}
}
{ =} { \begin{pmatrix}  xz \left(  y^2+z^2  \right) -xy^2z -xz \left(  x^2+y^2  \right)  \\ - x^2yz +yz \left(  x^2+z^2  \right) -yz \left(  x^2+y^2  \right) \\  x^2z^2 +y^2z^2 - \left(  x^2+y^2  \right)^2   \end{pmatrix}
}
{ =} { \left(  - x^2-y^2+z^2  \right) \begin{pmatrix}  xz \\ yz\\  x^2+y^2   \end{pmatrix}
}
{ =} { - \begin{pmatrix}  xz \\ yz\\  x^2+y^2   \end{pmatrix}
}
{ } {
}
}
{}
{}{.}
Der Eigenwert ist also
\mathl{\left(  x^2+y^2+z^2  \right)^{- { \frac{   3  }{  2 } } }}{}
\zusatzklammer {dies ist auch schon aus der beschreibenden Dreiecksmatrix ablesbar} {} {.}

}

__NOEDITSECTION__

\inputfaktbeweis
{Hyperfläche/Weingartenabbildung/Zweite Ableitung/Skalarprodukt/Fakt}
{Lemma}
{}
{

\faktsituation {Es sei

\mavergleichskette
{\vergleichskette
{ W
}
{ \subseteq }{ \R^n
}
{  }{
}
{    }{
}
{   }{
}
}
{}{}{}
\definitionsverweis {offen}{}{,}
\maabb {h} { W } { \R
} {}
eine zweifach
\definitionsverweis {stetig differenzierbare Funktion}{}{}
und

\mavergleichskette
{\vergleichskette
{ Y
}
{ = }{ h^{-1}(c)
}
{  }{
}
{    }{
}
{   }{
}
}
{}{}{}
die
\definitionsverweis {Faser}{}{}
zu

\mavergleichskette
{\vergleichskette
{ c
}
{ \in }{ \R
}
{  }{
}
{    }{
}
{   }{
}
}
{}{}{,}
wobei $h$ in jedem Punkt von $Y$
\definitionsverweis {regulär}{}{}
sei. Es sei $N$ ein
\definitionsverweis {Einheitsnormalenfeld}{}{}
und es sei

\mavergleichskette
{\vergleichskette
{ P
}
{ \in }{ Y
}
{  }{
}
{    }{
}
{   }{
}
}
{}{}{.}
Es sei $\gamma$ eine zweifach differenzierbare Realisierung auf $Y$ eines Tangentenvektors

\mavergleichskette
{\vergleichskette
{ v
}
{ \in }{ T_PY
}
{  }{
}
{    }{
}
{   }{
}
}
{}{}{.}}
\faktfolgerung {Dann ist

\mavergleichskettedisp
{\vergleichskette
{ \left\langle  \gamma^{\prime \prime} (0)  ,  N(P)  \right\rangle
}
{ =} { \left\langle  L_P(v) , v  \right\rangle
}
{ } {
}
{ } {
}
{ } {
}
}
{}{}{.}}
\faktzusatz {}
\faktzusatz {}

}
{

Es sei
\maabbeledisp {\gamma} { I } { Y
} { t } { \gamma(t)
} {,}
zweifach differenzierbar mit

\mavergleichskette
{\vergleichskette
{ \gamma(0)
}
{ = }{ P
}
{  }{
}
{    }{
}
{   }{
}
}
{}{}{}
und

\mavergleichskette
{\vergleichskette
{ \gamma'(0)
}
{ = }{ v
}
{ \in }{ T_PY
}
{    }{
}
{   }{
}
}
{}{}{.}
Es ist

\mavergleichskettedisp
{\vergleichskette
{ \left\langle  \gamma'(t)  ,  N(\gamma(t))   \right\rangle
}
{ =} { 0
}
{ } {
}
{ } {
}
{ } {
}
}
{}{}{}
für alle

\mavergleichskette
{\vergleichskette
{ t
}
{ \in }{ I
}
{  }{
}
{    }{
}
{   }{
}
}
{}{}{}
und daher ist mit
[[Skalarprodukt von Abbildungen/Richtungsableitung/Aufgabe|Aufgabe 43.14 (Analysis (Osnabrück 2021-2023))]]

\mavergleichskettealign
{\vergleichskettealign
{ 0
}
{ =} { \left\langle  \gamma'(t)  ,  N(\gamma(t))   \right\rangle' (0)
}
{ =} { \left\langle  \gamma^{\prime \prime}(0)  ,  N(\gamma(0))  \right\rangle + \left\langle  \gamma'(0)  ,  \left(  D_{ \gamma'(0) }  N   \right) \left(   \gamma(0)   \right)   \right\rangle
}
{ =} { \left\langle  \gamma^{\prime \prime}(0) , N(P)  \right\rangle - \left\langle  L_P(v)  ,  v   \right\rangle
}
{ } {
}
}
{}
{}{.}

}

In der vorstehenden Aussage ist keine zusätzliche Festlegung über die Orientierung nötig, da auf beiden Seiten das Einheitsnormalenfeld, rechts via die Weingartenabbildung, eingeht. Man beachte ferner, dass $\gamma$ keine konstante Geschwindigkeit besitzen muss, es also viele Realisierungen gibt und die Beschleunigung $\gamma^{\prime \prime} (0)$ nicht festgelegt ist. Der Ausdruck
\mathl{\left\langle  \gamma^{\prime \prime}  ,   \right\rangle(0)}{,} der die normale Komponente der Beschleunigung beschreibt, ändert sich aber nicht, da sich $\gamma$ auf der Hyperfläche bewegt.

__NOEDITSECTION__

\inputfaktbeweis
{Hyperfläche/Weingartenabbildung/Selbstadjungiert/Fakt}
{Satz}
{}
{

\faktsituation {Es sei

\mavergleichskette
{\vergleichskette
{ W
}
{ \subseteq }{ \R^n
}
{  }{
}
{    }{
}
{   }{
}
}
{}{}{}
\definitionsverweis {offen}{}{,}
\maabb {h} { W } { \R
} {}
eine zweifach
\definitionsverweis {stetig differenzierbare Funktion}{}{}
und

\mavergleichskette
{\vergleichskette
{ Y
}
{ = }{ h^{-1}(c)
}
{  }{
}
{    }{
}
{   }{
}
}
{}{}{}
die Faser zu

\mavergleichskette
{\vergleichskette
{ c
}
{ \in }{ \R
}
{  }{
}
{    }{
}
{   }{
}
}
{}{}{,}
wobei $h$ in jedem Punkt von $Y$
\definitionsverweis {regulär}{}{}
sei.}
\faktfolgerung {Dann ist die
\definitionsverweis {Weingartenabbildung}{}{}
$L_P$ in jedem Punkt

\mavergleichskette
{\vergleichskette
{ P
}
{ \in }{ Y
}
{  }{
}
{    }{
}
{   }{
}
}
{}{}{}
\definitionsverweis {selbstadjungiert}{}{.}}
\faktzusatz {}
\faktzusatz {}

}
{

Für Vektoren

\mavergleichskette
{\vergleichskette
{ v,w
}
{ \in }{ T_PY
}
{  }{
}
{    }{
}
{   }{
}
}
{}{}{}
ist

\mavergleichskettedisp
{\vergleichskette
{ \left\langle  v  ,  L_P(w)  \right\rangle
}
{ =} { \left\langle L_P(v) ,  w  \right\rangle
}
{ } {
}
{ } {
}
{ } {
}
}
{}{}{}
zu zeigen. Mit

\mavergleichskette
{\vergleichskette
{ N(P)
}
{ = }{ { \frac{
\operatorname{Grad} \,  h  (  P )  }{  \Vert {
\operatorname{Grad} \,  h  (  P ) } \Vert  } }
}
{  }{
}
{    }{
}
{   }{
}
}
{}{}{}
ist gemäß
[[Totale Differenzierbarkeit/R/Produktregel/Fakt|Lemma 45.10 (Analysis (Osnabrück 2021-2023))]]

\mavergleichskettealignhandlinks
{\vergleichskettealignhandlinks
{ \left\langle  v  ,  L_P(w)  \right\rangle
}
{ =} { -\left\langle  v  ,  \left(  D_{w}  N   \right) \left(   P   \right)   \right\rangle
}
{ =} { -\left\langle  v  ,  \left(  D_{w}  { \frac{
\operatorname{Grad} \,  h  }{  \Vert {
\operatorname{Grad} \,  h } \Vert  } }   \right) \left(   P   \right)     \right\rangle
}
{ =} { -\left\langle  v  ,  \left(  D_{w}  { \frac{   1  }{  \Vert {
\operatorname{Grad} \,  h } \Vert  } }   \right) \left(   P   \right) \cdot
\operatorname{Grad} \,  h  (  P ) + { \frac{   1  }{  \Vert {
\operatorname{Grad} \,  h  (  P ) } \Vert  } } \cdot  \left(  D_{w}
\operatorname{Grad} \,  h   \right) \left(   P   \right)   \right\rangle
}
{ =} { - { \frac{   1  }{  \Vert {
\operatorname{Grad} \,  h  (  P ) } \Vert  } } \left\langle  v  ,  \left(  D_{w}
\operatorname{Grad} \,  h   \right) \left(   P   \right)    \right\rangle
}
}
{}
{}{,}
da der erste Summand senkrecht auf dem Tangentialvektor $v$ steht. Mit Koordinatenfunktionen ist

\mavergleichskettealignhandlinks
{\vergleichskettealignhandlinks
{ \left(  D_{w}
\operatorname{Grad} \,  h   \right) \left(   P   \right)
}
{ =} { \left(  D_{w}  \left(  { \frac{   \partial   h   }{  \partial   x_1   } }   , \,   \ldots  , \,   { \frac{   \partial   h   }{  \partial   x_n   } }                 \right)    \right) \left(   P   \right)
}
{ =} { \sum_{j = 1}^n  w_j  { \frac{   \partial    }{  \partial   x_j   } }  \left(  { \frac{   \partial   h   }{  \partial   x_1   } }   , \,   \ldots  , \,   { \frac{   \partial   h   }{  \partial   x_n   } }                 \right) (P)
}
{ =} { \left(  \sum_{j = 1}^n  w_j  { \frac{   \partial    }{  \partial   x_j   } }  { \frac{   \partial   h   }{  \partial   x_1   } } (P)   , \,   \ldots  , \,   \sum_{j = 1}^n  w_j  { \frac{   \partial    }{  \partial   x_j   } }  { \frac{   \partial   h   }{  \partial   x_n   } } (P)                 \right)
}
{ } {
}
}
{}
{}{.}
Der obige Ausdruck ist somit gleich

\mavergleichskettealign
{\vergleichskettealign
{ \left\langle  v  ,  L_P(w)  \right\rangle
}
{ =} { - { \frac{   1  }{  \Vert {
\operatorname{Grad} \,  h  (  P ) } \Vert  } } \left\langle  v  ,  \left(  D_{w}
\operatorname{Grad} \,  h   \right) \left(   P   \right)   \right\rangle
}
{ =} { -  { \frac{   1  }{  \Vert {
\operatorname{Grad} \,  h  (  P ) } \Vert  } }  \sum_{i,j = 1}^n v_i  w_j  { \frac{   \partial    }{  \partial   x_j   } } { \frac{   \partial   h   }{  \partial   x_i   } }  (P)
}
{ } {
}
{ } {
}
}
{}
{}{.}
Nach
[[Differenzierbarkeit/Satz von Schwarz/Fakt|Satz 44.10 (Analysis (Osnabrück 2021-2023))]]
kann man die Reihenfolge der partiellen Ableitungen vertauschen, sodass man auch die Rollen von
\mathkor {} {v} {und} {w} {}
vertauschen kann.

}

__NOEDITSECTION__

\inputfaktbeweis
{Hyperfläche/Weingartenabbildung/Selbstadjungiert/Eigenwerte/Fakt}
{Korollar}
{}
{

\faktsituation {Es sei

\mavergleichskette
{\vergleichskette
{ W
}
{ \subseteq }{ \R^n
}
{  }{
}
{    }{
}
{   }{
}
}
{}{}{}
\definitionsverweis {offen}{}{,}
\maabb {h} { W } { \R
} {}
eine zweifach
\definitionsverweis {stetig differenzierbare Funktion}{}{}
und

\mavergleichskette
{\vergleichskette
{ Y
}
{ = }{ h^{-1}(c)
}
{  }{
}
{    }{
}
{   }{
}
}
{}{}{}
die Faser zu

\mavergleichskette
{\vergleichskette
{ c
}
{ \in }{ \R
}
{  }{
}
{    }{
}
{   }{
}
}
{}{}{,}
wobei $h$ in jedem Punkt von $Y$
\definitionsverweis {regulär}{}{}
sei.}
\faktfolgerung {Dann ist die
\definitionsverweis {Weingartenabbildung}{}{}
$L_P$ in jedem Punkt

\mavergleichskette
{\vergleichskette
{ P
}
{ \in }{ Y
}
{  }{
}
{    }{
}
{   }{
}
}
{}{}{}
\definitionsverweis {diagonalisierbar}{}{}
mit reellen
\definitionsverweis {Eigenwerten}{}{}
und zueinander orthogonalen Eigenräumen.}
\faktzusatz {}
\faktzusatz {}

}
{

Dies folgt aus
[[Hyperfläche/Weingartenabbildung/Selbstadjungiert/Fakt|Satz 4.7]]
und
[[Vektorraum mit Skalarprodukt/Endomorphismus/Selbstadjungiert/Spektralsatz/Fakt|Satz 41.11 (Lineare Algebra (Osnabrück 2024-2025))]].

}

__NOEDITSECTION__

\inputfaktbeweis
{Differenzierbare Funktion/Graph/Hyperfläche/Weingartenabbildung/Fakt}
{Lemma}
{}
{

\faktsituation {Es sei

\mavergleichskette
{\vergleichskette
{ V
}
{ \subseteq }{ \R^{n-1}
}
{  }{
}
{    }{
}
{   }{
}
}
{}{}{}
eine
\definitionsverweis {offene Menge}{}{}
und sei

\mavergleichskette
{\vergleichskette
{ Y
}
{ \subseteq }{ V \times \R
}
{  }{
}
{    }{
}
{   }{
}
}
{}{}{}
der
\definitionsverweis {Graph}{}{}
zu einer zweifach
\definitionsverweis {stetig differenzierbaren}{}{}
Funktion
\maabbdisp {f} { V } {\R
} {.}}
\faktfolgerung {Dann ist die
\definitionsverweis {Weingartenabbildung}{}{}
$L_P$ in einem Punkt

\mavergleichskette
{\vergleichskette
{ P
}
{ = }{ (x_1 , \ldots , x_{n-1}, f(x_1 , \ldots , x_{n-1}))
}
{ \in }{ Y
}
{    }{
}
{   }{
}
}
{}{}{}
\zusatzklammer {zu dem nach oben gerichteten Einheitsnormalenfeld} {} {}
durch
\mathl{{ \frac{   1  }{  \sqrt{ 1+ \left(  { \frac{   \partial   f   }{  \partial   x_1   } }  \right)^2 + \cdots + \left(  { \frac{   \partial   f   }{  \partial   x_{n-1}   } }  \right)^2 }  } } \operatorname{Hess}_{   } \,  f}{} gegeben, wobei $\operatorname{Hess}_{   } \,  f$ die
\definitionsverweis {Hesse-Matrix}{}{}
zu $f$ bezeichnet und wenn man Grundvektoren

\mavergleichskette
{\vergleichskette
{ v
}
{ \in }{ \R^{n-1}
}
{  }{
}
{    }{
}
{   }{
}
}
{}{}{}
mit den Tangentialvektoren aus $T_PY$ im Sinne von
[[Funktion/n-1 Variablen/Graph als Hyperfläche/Tangentialraum/Beispiel|Beispiel 1.2]]
identifiziert.}
\faktzusatz {}
\faktzusatz {}

}
{

Wir knüpfen an
[[Funktion/n-1 Variablen/Graph als Hyperfläche/Tangentialraum/Beispiel|Beispiel 1.2]]
an. Das nach oben gerichtete Einheitsnormalenfeld ist durch

\mavergleichskettedisphandlinks
{\vergleichskettedisphandlinks
{ N(x_1  , \ldots , x_{n-1}, f(x_1 , \ldots , x_{n-1}))
}
{ =} { { \frac{   1  }{  \sqrt{ 1+ \left(  { \frac{   \partial   f   }{  \partial   x_1   } }  \right)^2 + \cdots + \left(  { \frac{   \partial   f   }{  \partial   x_{n-1}   } }  \right)^2 }  } } \begin{pmatrix}  -{ \frac{   \partial   f   }{  \partial   x_{1}   } }  \\\vdots\\  -{ \frac{   \partial   f   }{  \partial   x_{n-1}   } }\\ 1   \end{pmatrix}
}
{ } {
}
{ } {
}
{ } {
}
}
{}{}{}
gegeben. Wir betrachten den Weg

\mavergleichskettedisp
{\vergleichskette
{ \gamma(t)
}
{ =} { \begin{pmatrix}  x_1+tv_1  \\ \vdots \\  x_{n-1} +tv_{n-1}\\f \begin{pmatrix}  x_1+tv_1  \\ \vdots \\  x_{n-1} +tv_{n-1}   \end{pmatrix}   \end{pmatrix}
}
{ } {
}
{ } {
}
{ } {
}
}
{}{}{}
auf dem Graphen zum Grundvektor $v$. Die zweite Ableitung davon ist

\mavergleichskettedisp
{\vergleichskette
{ \gamma^{\prime \prime} (0)
}
{ =} { \begin{pmatrix}  0  \\ \vdots \\  0\\ \sum_{1 \leq i, j \leq n-1}  { \frac{   \partial    }{  \partial   x_i   } }  { \frac{   \partial   f   }{  \partial   x_j   } } v_iv_j   \end{pmatrix}
}
{ } {
}
{ } {
}
{ } {
}
}
{}{}{.}
Nach
[[Hyperfläche/Weingartenabbildung/Zweite Ableitung/Skalarprodukt/Fakt|Lemma 4.6]]
ist

\mavergleichskettealigndrucklinks
{\vergleichskettealigndrucklinks
{ \left\langle  L_P( \gamma'(0) ) ,  \gamma'(0)  \right\rangle
}
{ =} { \left\langle  \gamma^{\prime \prime} (0) ,  N(x_1 , \ldots , x_{n-1}, f(x_1  , \ldots , x_{n-1}))   \right\rangle
}
{ =} { { \frac{   1  }{  \sqrt{ 1 + \left(  { \frac{   \partial   f   }{  \partial   x_1   } }  \right)^2 + \cdots + \left(  { \frac{   \partial   f  }{  \partial   x_{n-1}   } }  \right)^2 }  } } \left\langle  \begin{pmatrix}  0  \\\vdots\\  0\\ \sum_{ 1 \leq i, j \leq n-1 }  { \frac{   \partial    }{  \partial   x_i   } }  { \frac{   \partial   f   }{  \partial   x_j   } } v_iv_j   \end{pmatrix}  ,  \begin{pmatrix}  -{ \frac{   \partial   f   }{  \partial   x_{1}   } }  \\\vdots\\  -{ \frac{   \partial   f   }{  \partial   x_{n-1}   } }\\ 1   \end{pmatrix}   \right\rangle
}
{ =} { { \frac{   \sum_{ 1 \leq i, j \leq n-1 } { \frac{   \partial    }{  \partial   x_i   } }  { \frac{   \partial   f   }{  \partial   x_j   } } v_i v_j     }{  \sqrt{ 1+  \left(  { \frac{   \partial   f   }{  \partial   x_1   } }  \right)^2 + \cdots + \left(  { \frac{   \partial   f   }{  \partial   x_{n-1}   } }  \right)^2 }  } }
}
{ =} { { \frac{   1  }{  \sqrt{ 1+ \left(  { \frac{   \partial   f   }{  \partial   x_1   } }  \right)^2 + \cdots + \left(  { \frac{   \partial   f   }{  \partial   x_{n-1}   } }  \right)^2 }  } }  \left( v_1   , \,   \ldots  , \,   v_{n-1}                 \right)  \operatorname{Hess}_{   } \,  f  \begin{pmatrix} v_1  \\ \vdots \\  v_{n-1}   \end{pmatrix}
}
}
{}{}{.}
Das bedeutet, dass die nach
[[Hyperfläche/Weingartenabbildung/Selbstadjungiert/Fakt|Satz 4.7]]
symmetrische
\definitionsverweis {Bilinearform}{}{}

\mathl{\left\langle  L_P(-) ,  -  \right\rangle}{} im Tangentialraum $T_PY$ mit der durch die
\zusatzklammer {durch den Vorfaktor} {} {}
skalierte Hessematrix gegebenen Bilinearform auf $\R^{n-1}$ übereinstimmt, wenn vorne und hinten der gleiche Vektor eingesetzt wird. Nach
[[Bilinearform/Symmetrisch/Polarisationsformel/Fakt|Lemma 38.10 (Lineare Algebra (Osnabrück 2024-2025))]]
stimmen dann generell die Bilinearformen über ein. Dann stimmen auch die linearen Abbildungen $L_P$ und die durch die Hessematrix gegebene lineare Abbildung überein.

}

 [[Kategorie:Latexseite]]
